\chapter{RAG Pipeline \& Ablation Benchmark}
\label{ch:ablation}

\section{Pilot Run: Notebook 04}

A pilot evaluation was run using Qwen3-4B (standard, non-thinking variant) with Qwen3-Embedding-0.6B on 100 MedQA questions. Retrieval used a \gls{faiss} \gls{ifip} \cite{johnson_billionscale_2019} over 36,723 textbook chunks. Results: \gls{rag} accuracy = 38\%, baseline = 42\%, delta = $-4$\gls{pp}. RAG helped 6 questions and hurt 8. Mean retrieval score = 0.648. This result prompted the hypothesis that the standard \gls{llm} variant cannot effectively use retrieved context. 

\section{Full Ablation Grid: Notebook 07}
A six-run ablation grid was conducted with varying embedding models, \gls{llm} variants, and \gls{rag} conditions. All runs evaluated on MedQA questions; runs 3-6 used 100 questions each. See Table \ref{tab:ablation} in Appendix for Grid Results, and Figure \ref{fig:ablation} for visual. For reference, \gls{amg-rag} achieves 74.1\% F1 on MedQA, powered by GPT-4o-mini, which is a paid model with $\sim$8B parameters, integrated with a dynamic PubMed knowledge graph \cite{rezaei_agentic_2025}. 

\section{Interpretation}

The ablation reveals three findings. First, the standard Qwen3-4B cannot use retrieved context effectively: \gls{rag} hurts by $-4$\gls{pp} regardless of embedding. Second, switching to the Thinking variant alone is insufficient: with Qwen3-Embedding, \gls{rag} adds only 1\gls{pp}. Third, the combination of a biomedical-grade embedding with a reasoning-capable \gls{llm} (Thinking-2507) is required for \gls{rag} to help: together they deliver +11\gls{pp}.

This answers the research question affirmatively but conditionally: \gls{rag} improves accuracy, but only when the embedding is biomedical-grade and the \gls{llm} is capable of reasoning over retrieved passages. The 30\gls{pp} gap to \gls{amg-rag} (74.1\%) is attributable to three factors: \gls{llm} parameter scale (4B vs. GPT-4o-mini's $\sim$8B), static textbook retrieval versus \gls{amg-rag}'s dynamic live PubMed knowledge graph construction, and the absence of fine-tuning on the MedQA retrieval task specifically. Both systems evaluate on MedQA, so dataset difficulty is not a differentiating factor.